\documentclass[11pt,a4paper]{article}
\usepackage[utf8]{inputenc}
\usepackage[T1]{fontenc}
\usepackage{amsmath,amssymb}
\usepackage{graphicx}
\usepackage{booktabs}
\usepackage{siunitx}
\usepackage{geometry}
\geometry{margin=1in}
\usepackage{pythontex}
\usepackage{hyperref}
\usepackage{float}

\title{Weber-Fechner Law\\JND}
\author{Psychophysics Research Group}
\date{\today}

\begin{document}
\maketitle

\begin{abstract}
This document presents a comprehensive computational analysis of classical psychophysical laws governing the relationship between physical stimulus intensity and perceived sensation magnitude. We implement Weber's Law for just noticeable differences, Fechner's logarithmic law, Stevens' power law, psychometric functions, and the method of constant stimuli. The analysis demonstrates how different sensory modalities exhibit distinct mathematical relationships between objective stimulus intensity and subjective sensation, with implications for sensory system design and perceptual modeling.
\end{abstract}

\section{Introduction}

The Weber-Fechner Law represents one of the foundational principles in psychophysics, describing the quantitative relationship between physical stimulus intensity and perceived sensation. Weber's Law (1834) states that the just noticeable difference (JND) in stimulus intensity is proportional to the base intensity: $\Delta I = k \cdot I$, where $k$ is the Weber fraction. Fechner (1860) extended this to propose that sensation magnitude grows logarithmically with stimulus intensity: $S = c \log(I/I_0)$. Stevens (1957) later proposed an alternative power law relationship: $S = k I^n$, where the exponent $n$ varies across sensory modalities. This analysis implements all three formulations alongside psychometric functions and threshold estimation methods, providing a complete computational framework for psychophysical analysis.

\begin{pycode}

import numpy as np
import matplotlib.pyplot as plt
from scipy import stats, optimize, integrate
from scipy.stats import norm
plt.rcParams['text.usetex'] = True
plt.rcParams['font.family'] = 'serif'

\end{pycode}

\section{Weber's Law: Just Noticeable Difference}

Weber's Law states that the just noticeable difference (JND) is proportional to the stimulus intensity. This fundamental principle applies across diverse sensory modalities including vision, audition, and touch.

\begin{pycode}
# Weber's Law: delta_I = k * I
I = np.linspace(1, 100, 100)  # stimulus intensity
k_vision = 0.08  # Weber fraction for brightness
k_weight = 0.02  # Weber fraction for weight perception
k_audition = 0.10  # Weber fraction for loudness

delta_I_vision = k_vision * I
delta_I_weight = k_weight * I
delta_I_audition = k_audition * I

fig, (ax1, ax2) = plt.subplots(1, 2, figsize=(12, 5))

# Plot 1: JND vs Intensity for different modalities
ax1.plot(I, delta_I_vision, 'b-', linewidth=2, label='Vision (k=0.08)')
ax1.plot(I, delta_I_weight, 'r-', linewidth=2, label='Weight (k=0.02)')
ax1.plot(I, delta_I_audition, 'g-', linewidth=2, label='Audition (k=0.10)')
ax1.set_xlabel('Stimulus Intensity $I$', fontsize=12)
ax1.set_ylabel('Just Noticeable Difference $\\Delta I$', fontsize=12)
ax1.set_title("Weber's Law: JND vs Stimulus Intensity", fontsize=13)
ax1.legend(loc='upper left', fontsize=10)
ax1.grid(True, alpha=0.3)

# Plot 2: Weber fraction (constant ratio)
weber_ratio_vision = delta_I_vision / I
weber_ratio_weight = delta_I_weight / I
ax2.plot(I, weber_ratio_vision, 'b-', linewidth=2, label='Vision')
ax2.plot(I, weber_ratio_weight, 'r-', linewidth=2, label='Weight')
ax2.plot(I, weber_ratio_audition / I, 'g-', linewidth=2, label='Audition')
ax2.set_xlabel('Stimulus Intensity $I$', fontsize=12)
ax2.set_ylabel('Weber Fraction $\\Delta I / I$', fontsize=12)
ax2.set_title('Constant Weber Fraction Across Intensities', fontsize=13)
ax2.legend(loc='upper right', fontsize=10)
ax2.grid(True, alpha=0.3)
ax2.set_ylim([0, 0.15])

plt.tight_layout()
plt.savefig('weber_fechner_plot1.pdf', dpi=150, bbox_inches='tight')
plt.close()

# Store values for inline reference
vision_jnd_at_50 = k_vision * 50
weight_jnd_at_50 = k_weight * 50
\end{pycode}

\begin{figure}[H]
\centering
\includegraphics[width=0.95\textwidth]{weber_fechner_plot1.pdf}
\caption{Weber's Law demonstration showing linear relationship between just noticeable difference (JND) and stimulus intensity across three sensory modalities. Left panel shows absolute JND values increasing proportionally with intensity. Right panel confirms constant Weber fraction across intensities, validating Weber's principle. At intensity level 50, vision requires JND of \py{vision_jnd_at_50:.2f} units while weight perception requires only \py{weight_jnd_at_50:.2f} units due to different Weber fractions.}
\end{figure}

\section{Fechner's Law: Logarithmic Sensation}

Fechner proposed that sensation magnitude grows logarithmically with stimulus intensity: $S = c \log(I/I_0)$, where $I_0$ is the threshold intensity and $c$ is a scaling constant. This relationship emerges from integrating Weber's Law across intensity levels.

\begin{pycode}
# Fechner's Law: S = c * log(I / I_0)
I_fechner = np.linspace(1, 100, 100)
I_0 = 1.0  # threshold intensity
c = 10.0   # scaling constant

S_fechner = c * np.log(I_fechner / I_0)

# Compare with linear relationship
S_linear = 0.5 * I_fechner  # hypothetical linear sensation

fig, (ax1, ax2) = plt.subplots(1, 2, figsize=(12, 5))

# Plot 1: Fechner's logarithmic law
ax1.plot(I_fechner, S_fechner, 'b-', linewidth=2, label='Fechner (logarithmic)')
ax1.plot(I_fechner, S_linear, 'r--', linewidth=2, label='Linear (comparison)')
ax1.set_xlabel('Stimulus Intensity $I$', fontsize=12)
ax1.set_ylabel('Sensation Magnitude $S$', fontsize=12)
ax1.set_title("Fechner's Law: Logarithmic Sensation Growth", fontsize=13)
ax1.legend(loc='upper left', fontsize=10)
ax1.grid(True, alpha=0.3)

# Plot 2: Derivative showing diminishing returns
dS_dI = c / I_fechner  # derivative of Fechner's law
ax2.plot(I_fechner, dS_dI, 'g-', linewidth=2)
ax2.set_xlabel('Stimulus Intensity $I$', fontsize=12)
ax2.set_ylabel('Sensitivity $dS/dI$', fontsize=12)
ax2.set_title('Diminishing Sensitivity with Increasing Intensity', fontsize=13)
ax2.grid(True, alpha=0.3)

plt.tight_layout()
plt.savefig('weber_fechner_plot2.pdf', dpi=150, bbox_inches='tight')
plt.close()

# Store values for inline reference
sensation_at_10 = c * np.log(10 / I_0)
sensation_at_100 = c * np.log(100 / I_0)
\end{pycode}

\begin{figure}[H]
\centering
\includegraphics[width=0.95\textwidth]{weber_fechner_plot2.pdf}
\caption{Fechner's logarithmic law compared to linear sensation model. Left panel demonstrates compression of high-intensity stimuli: doubling intensity from 10 to 20 produces sensation increase from \py{sensation_at_10:.2f} to \py{c * np.log(20/I_0):.2f}, while doubling from 50 to 100 yields smaller perceptual difference. Right panel shows sensitivity ($dS/dI$) decreasing hyperbolically with intensity, explaining reduced discriminability at high intensities characteristic of sensory adaptation.}
\end{figure}

\section{Stevens' Power Law: Modality-Specific Exponents}

Stevens proposed that sensation magnitude follows a power law: $S = k I^n$, where the exponent $n$ characterizes different sensory modalities. This formulation better captures the diverse range of psychophysical relationships observed across senses.

\begin{pycode}
# Stevens' Power Law: S = k * I^n
I_stevens = np.linspace(1, 100, 100)
k = 1.0  # scaling constant

# Different exponents for different modalities
n_brightness = 0.33   # visual brightness (compressive)
n_loudness = 0.67     # auditory loudness (slightly compressive)
n_length = 1.0        # visual length (veridical)
n_electric = 3.5      # electric shock intensity (expansive)

S_brightness = k * I_stevens**n_brightness
S_loudness = k * I_stevens**n_loudness
S_length = k * I_stevens**n_length
S_electric = k * I_stevens**n_electric

fig, ax = plt.subplots(figsize=(10, 6))

ax.plot(I_stevens, S_brightness, 'b-', linewidth=2, label=f'Brightness (n={n_brightness})')
ax.plot(I_stevens, S_loudness, 'g-', linewidth=2, label=f'Loudness (n={n_loudness})')
ax.plot(I_stevens, S_length, 'orange', linewidth=2, label=f'Length (n={n_length})')
ax.plot(I_stevens, S_electric, 'r-', linewidth=2, label=f'Electric Shock (n={n_electric})')

ax.set_xlabel('Stimulus Intensity $I$', fontsize=12)
ax.set_ylabel('Sensation Magnitude $S$', fontsize=12)
ax.set_title("Stevens' Power Law: Modality-Specific Exponents", fontsize=13)
ax.legend(loc='upper left', fontsize=10)
ax.grid(True, alpha=0.3)
ax.set_ylim([0, 50])

plt.tight_layout()
plt.savefig('weber_fechner_plot3.pdf', dpi=150, bbox_inches='tight')
plt.close()

# Store values for inline reference
brightness_at_50 = k * 50**n_brightness
electric_at_50 = k * 50**n_electric
\end{pycode}

\begin{figure}[H]
\centering
\includegraphics[width=0.85\textwidth]{weber_fechner_plot3.pdf}
\caption{Stevens' power law demonstrates diverse exponent values across sensory modalities. Compressive functions ($n < 1$) characterize brightness and loudness, where high-intensity stimuli produce smaller-than-proportional sensation increases. Veridical perception ($n = 1$) occurs for visual length estimation. Expansive functions ($n > 1$) describe pain and electric shock, where sensation magnitude grows dramatically: at intensity 50, brightness sensation reaches \py{brightness_at_50:.2f} while electric shock sensation reaches \py{electric_at_50:.2f}, reflecting protective amplification mechanisms.}
\end{figure}

\section{Psychometric Function and Detection Thresholds}

The psychometric function describes the probability of stimulus detection as a function of intensity. It typically follows a sigmoid shape, characterized by the detection threshold (50\% detection point) and slope (sensitivity).

\begin{pycode}
# Psychometric function: cumulative Gaussian
I_psychometric = np.linspace(0, 100, 200)
threshold = 50  # 50% detection point
slope = 10      # smaller = steeper function (higher sensitivity)

# Cumulative Gaussian (standard psychometric function)
p_detect = norm.cdf(I_psychometric, loc=threshold, scale=slope)

# Alternative: Weibull function (common in psychophysics)
from scipy.special import erf
alpha = 50  # threshold parameter
beta = 2    # shape parameter
p_detect_weibull = 1 - np.exp(-(I_psychometric/alpha)**beta)

fig, (ax1, ax2) = plt.subplots(1, 2, figsize=(12, 5))

# Plot 1: Standard psychometric function
ax1.plot(I_psychometric, p_detect, 'b-', linewidth=2)
ax1.axhline(y=0.5, color='r', linestyle='--', linewidth=1, label='Threshold (50%)')
ax1.axvline(x=threshold, color='r', linestyle='--', linewidth=1)
ax1.axhline(y=0.75, color='g', linestyle=':', linewidth=1, label='75% detection')
ax1.set_xlabel('Stimulus Intensity $I$', fontsize=12)
ax1.set_ylabel('Detection Probability $P(detect)$', fontsize=12)
ax1.set_title('Psychometric Function (Cumulative Gaussian)', fontsize=13)
ax1.legend(loc='lower right', fontsize=10)
ax1.grid(True, alpha=0.3)
ax1.set_ylim([0, 1])

# Plot 2: Comparison of different slopes
slopes = [5, 10, 20]
for s in slopes:
    p = norm.cdf(I_psychometric, loc=threshold, scale=s)
    ax2.plot(I_psychometric, p, linewidth=2, label=f'$\\sigma$ = {s}')
ax2.set_xlabel('Stimulus Intensity $I$', fontsize=12)
ax2.set_ylabel('Detection Probability $P(detect)$', fontsize=12)
ax2.set_title('Effect of Slope on Sensitivity', fontsize=13)
ax2.legend(loc='lower right', fontsize=10)
ax2.grid(True, alpha=0.3)
ax2.set_ylim([0, 1])

plt.tight_layout()
plt.savefig('weber_fechner_plot4.pdf', dpi=150, bbox_inches='tight')
plt.close()

# Store values for inline reference
prob_at_threshold = norm.cdf(threshold, loc=threshold, scale=slope)
intensity_75_percent = norm.ppf(0.75, loc=threshold, scale=slope)
\end{pycode}

\begin{figure}[H]
\centering
\includegraphics[width=0.95\textwidth]{weber_fechner_plot4.pdf}
\caption{Psychometric functions model detection probability as sigmoid function of stimulus intensity. Left panel shows standard cumulative Gaussian with threshold at intensity \py{threshold:.1f} where detection probability is \py{prob_at_threshold:.2f}. Seventy-five percent detection occurs at intensity \py{intensity_75_percent:.2f}. Right panel demonstrates slope effects: smaller slope values produce steeper functions indicating higher observer sensitivity and more precise threshold estimation. Shallow slopes suggest higher measurement noise or greater inter-trial variability.}
\end{figure}

\section{Method of Constant Stimuli}

The method of constant stimuli is a classical psychophysical technique where a fixed set of stimulus intensities is presented repeatedly in random order, and the observer's detection rate is measured at each level. This data is fit to a psychometric function to estimate the detection threshold.

\begin{pycode}
# Method of Constant Stimuli
np.random.seed(42)

# Fixed stimulus intensities tested
intensities = np.array([30, 40, 50, 60, 70, 80])

# Simulated detection rates (proportion detected out of n trials)
# True threshold = 55, so lower intensities have lower detection rates
n_trials = 100
true_threshold = 55
true_slope = 8

# Generate realistic detection data
detection_probs = norm.cdf(intensities, loc=true_threshold, scale=true_slope)
detections = np.array([np.random.binomial(n_trials, p) / n_trials for p in detection_probs])

# Add some experimental noise
detections = np.clip(detections + np.random.normal(0, 0.05, len(detections)), 0, 1)

fig, (ax1, ax2) = plt.subplots(1, 2, figsize=(12, 5))

# Plot 1: Raw data with fitted psychometric function
ax1.scatter(intensities, detections, s=100, color='red', label='Observed data', zorder=3)

# Fit psychometric function to data
from scipy.optimize import curve_fit
def psychometric(x, threshold, slope):
    return norm.cdf(x, loc=threshold, scale=slope)

params, _ = curve_fit(psychometric, intensities, detections, p0=[50, 10])
fitted_threshold, fitted_slope = params

I_fit = np.linspace(20, 90, 200)
p_fit = psychometric(I_fit, fitted_threshold, fitted_slope)

ax1.plot(I_fit, p_fit, 'b-', linewidth=2, label=f'Fitted function (threshold={fitted_threshold:.1f})')
ax1.axhline(y=0.5, color='gray', linestyle='--', linewidth=1, alpha=0.5)
ax1.axvline(x=fitted_threshold, color='gray', linestyle='--', linewidth=1, alpha=0.5)
ax1.set_xlabel('Stimulus Intensity', fontsize=12)
ax1.set_ylabel('Proportion Detected', fontsize=12)
ax1.set_title('Method of Constant Stimuli: Data and Fit', fontsize=13)
ax1.legend(loc='lower right', fontsize=10)
ax1.grid(True, alpha=0.3)
ax1.set_ylim([0, 1])

# Plot 2: Residuals
residuals = detections - psychometric(intensities, fitted_threshold, fitted_slope)
ax2.scatter(intensities, residuals, s=100, color='red')
ax2.axhline(y=0, color='black', linestyle='-', linewidth=1)
ax2.set_xlabel('Stimulus Intensity', fontsize=12)
ax2.set_ylabel('Residual (Observed - Fitted)', fontsize=12)
ax2.set_title('Fit Residuals', fontsize=13)
ax2.grid(True, alpha=0.3)

plt.tight_layout()
plt.savefig('weber_fechner_plot5.pdf', dpi=150, bbox_inches='tight')
plt.close()

# Store values for inline reference
threshold_estimate = fitted_threshold
slope_estimate = fitted_slope
rmse = np.sqrt(np.mean(residuals**2))
\end{pycode}

\begin{figure}[H]
\centering
\includegraphics[width=0.95\textwidth]{weber_fechner_plot5.pdf}
\caption{Method of constant stimuli demonstrates classical threshold estimation. Left panel shows six fixed intensity levels presented 100 times each, with detection proportion plotted against intensity. Fitted psychometric function yields estimated threshold of \py{threshold_estimate:.2f} with slope parameter \py{slope_estimate:.2f}. Right panel displays residuals with RMSE of \py{rmse:.3f}, confirming good fit quality. This method provides robust threshold estimates but requires extensive data collection across multiple intensity levels.}
\end{figure}

\section{Signal Detection Theory: ROC Analysis}

Signal Detection Theory (SDT) extends psychophysics by separating sensitivity (d-prime) from response bias. The Receiver Operating Characteristic (ROC) curve plots hit rate against false alarm rate across different decision criteria.

\begin{pycode}
# Signal Detection Theory: ROC curve
np.random.seed(42)

# Generate signal and noise distributions
n_samples = 10000
noise_dist = np.random.normal(0, 1, n_samples)
signal_dist = np.random.normal(1.5, 1, n_samples)  # d' = 1.5

# Calculate ROC curve by varying decision criterion
criteria = np.linspace(-2, 4, 100)
hit_rates = []
false_alarm_rates = []

for criterion in criteria:
    hits = np.sum(signal_dist >= criterion) / n_samples
    false_alarms = np.sum(noise_dist >= criterion) / n_samples
    hit_rates.append(hits)
    false_alarm_rates.append(false_alarms)

hit_rates = np.array(hit_rates)
false_alarm_rates = np.array(false_alarm_rates)

# Calculate d-prime and area under ROC curve
d_prime = 1.5  # known from distributions
auc = np.trapz(hit_rates, false_alarm_rates)

fig, (ax1, ax2) = plt.subplots(1, 2, figsize=(12, 5))

# Plot 1: Signal and noise distributions
x_range = np.linspace(-4, 6, 200)
ax1.plot(x_range, norm.pdf(x_range, 0, 1), 'b-', linewidth=2, label='Noise')
ax1.plot(x_range, norm.pdf(x_range, 1.5, 1), 'r-', linewidth=2, label='Signal')
ax1.axvline(x=0.75, color='g', linestyle='--', linewidth=1.5, label='Criterion')
ax1.fill_between(x_range, 0, norm.pdf(x_range, 0, 1), where=(x_range >= 0.75), alpha=0.3, color='blue', label='False Alarms')
ax1.fill_between(x_range, 0, norm.pdf(x_range, 1.5, 1), where=(x_range >= 0.75), alpha=0.3, color='red', label='Hits')
ax1.set_xlabel('Evidence', fontsize=12)
ax1.set_ylabel('Probability Density', fontsize=12)
ax1.set_title(f'Signal and Noise Distributions ($d\'$ = {d_prime})', fontsize=13)
ax1.legend(loc='upper right', fontsize=9)
ax1.grid(True, alpha=0.3)

# Plot 2: ROC curve
ax2.plot(false_alarm_rates, hit_rates, 'b-', linewidth=2, label=f'ROC curve (AUC={auc:.3f})')
ax2.plot([0, 1], [0, 1], 'k--', linewidth=1, label='Chance performance')
ax2.set_xlabel('False Alarm Rate', fontsize=12)
ax2.set_ylabel('Hit Rate', fontsize=12)
ax2.set_title('Receiver Operating Characteristic (ROC)', fontsize=13)
ax2.legend(loc='lower right', fontsize=10)
ax2.grid(True, alpha=0.3)
ax2.set_xlim([0, 1])
ax2.set_ylim([0, 1])

plt.tight_layout()
plt.savefig('weber_fechner_plot6.pdf', dpi=150, bbox_inches='tight')
plt.close()

# Store values for inline reference
auc_value = auc
dprime_value = d_prime
\end{pycode}

\begin{figure}[H]
\centering
\includegraphics[width=0.95\textwidth]{weber_fechner_plot6.pdf}
\caption{Signal Detection Theory analysis separates perceptual sensitivity from decision bias. Left panel shows overlapping signal and noise distributions with $d' = $ \py{dprime_value:.2f}, where decision criterion (green line) determines trade-off between hit rate (red shaded area) and false alarm rate (blue shaded area). Right panel displays ROC curve with area under curve of \py{auc_value:.3f}, indicating detection performance substantially above chance. Moving criterion leftward increases both hits and false alarms; rightward movement decreases both, but sensitivity remains constant.}
\end{figure}

\section{Results Summary}

\begin{pycode}
results = [
    ['Weber Fraction (Vision)', f'{k_vision:.3f}'],
    ['Weber Fraction (Weight)', f'{k_weight:.3f}'],
    ['Weber Fraction (Audition)', f'{k_audition:.3f}'],
    ['Fechner Scaling Constant', f'{c:.2f}'],
    ['Stevens Exponent (Brightness)', f'{n_brightness:.2f}'],
    ['Stevens Exponent (Loudness)', f'{n_loudness:.2f}'],
    ['Stevens Exponent (Electric Shock)', f'{n_electric:.2f}'],
    ['Detection Threshold (MCS)', f'{threshold_estimate:.2f}'],
    ['Psychometric Slope', f'{slope_estimate:.2f}'],
    ['SDT Sensitivity (d-prime)', f'{dprime_value:.2f}'],
    ['ROC Area Under Curve', f'{auc_value:.3f}'],
]

print(r'\\begin{table}[H]')
print(r'\\centering')
print(r'\\caption{Summary of Psychophysical Parameters}')
print(r'\\begin{tabular}{@{}lc@{}}')
print(r'\\toprule')
print(r'Parameter & Value \\\\')
print(r'\\midrule')
for row in results:
    print(f"{row[0]} & {row[1]} \\\\\\\\")
print(r'\\bottomrule')
print(r'\\end{tabular}')
print(r'\\end{table}')
\end{pycode}

\section{Conclusions}

This comprehensive analysis demonstrates the evolution of psychophysical theory from Weber's linear just noticeable difference law through Fechner's logarithmic integration to Stevens' power law formulation. Weber's Law successfully predicts constant discrimination ratios across intensities, with Weber fractions of \py{k_vision:.3f} for vision and \py{k_weight:.3f} for weight perception. Fechner's logarithmic law captures sensory compression but fails to account for modality-specific differences in scaling. Stevens' power law provides superior flexibility through exponent parameters: compressive functions ($n < 1$) for brightness (\py{n_brightness:.2f}) and loudness (\py{n_loudness:.2f}), veridical scaling ($n = 1$) for length estimation, and expansive functions ($n > 1$) for pain and electric shock (\py{n_electric:.2f}).

The method of constant stimuli yielded threshold estimates of \py{threshold_estimate:.2f} with slope \py{slope_estimate:.2f}, demonstrating classical psychometric function fitting. Signal Detection Theory extends this framework by separating sensitivity ($d' = $ \py{dprime_value:.2f}) from response bias, with ROC analysis confirming robust detection performance (AUC = \py{auc_value:.3f}). These computational methods form the foundation of modern sensory psychophysics, enabling quantitative modeling of perception across diverse modalities and experimental paradigms. Future work should incorporate adaptive psychophysical methods, Bayesian parameter estimation, and neural network models of sensory processing.

\section*{References}

\begin{enumerate}
\item Weber, E.H. (1834). \textit{De Pulsu, Resorptione, Auditu et Tactu: Annotationes Anatomicae et Physiologicae}. Leipzig: Koehler.

\item Fechner, G.T. (1860). \textit{Elemente der Psychophysik} [Elements of Psychophysics]. Leipzig: Breitkopf und Härtel.

\item Stevens, S.S. (1957). On the psychophysical law. \textit{Psychological Review}, 64(3), 153-181.

\item Stevens, S.S. (1961). To honor Fechner and repeal his law. \textit{Science}, 133(3446), 80-86.

\item Green, D.M., \& Swets, J.A. (1966). \textit{Signal Detection Theory and Psychophysics}. New York: Wiley.

\item Gescheider, G.A. (1997). \textit{Psychophysics: The Fundamentals} (3rd ed.). Mahwah, NJ: Lawrence Erlbaum Associates.

\item Kingdom, F.A.A., \& Prins, N. (2016). \textit{Psychophysics: A Practical Introduction} (2nd ed.). London: Academic Press.

\item Macmillan, N.A., \& Creelman, C.D. (2005). \textit{Detection Theory: A User's Guide} (2nd ed.). Cambridge: Cambridge University Press.

\item Treutwein, B. (1995). Adaptive psychophysical procedures. \textit{Vision Research}, 35(17), 2503-2522.

\item Wichmann, F.A., \& Hill, N.J. (2001). The psychometric function: I. Fitting, sampling, and goodness of fit. \textit{Perception \& Psychophysics}, 63(8), 1293-1313.

\item Pelli, D.G., \& Farell, B. (1999). Why use noise? \textit{Journal of the Optical Society of America A}, 16(3), 647-653.

\item Kontsevich, L.L., \& Tyler, C.W. (1999). Bayesian adaptive estimation of psychometric slope and threshold. \textit{Vision Research}, 39(16), 2729-2737.

\item Norwich, K.H. (1993). \textit{Information, Sensation, and Perception}. San Diego: Academic Press.

\item Laming, D. (1986). \textit{Sensory Analysis}. London: Academic Press.

\item Wickens, T.D. (2002). \textit{Elementary Signal Detection Theory}. Oxford: Oxford University Press.

\item Knoblauch, K., \& Maloney, L.T. (2012). \textit{Modeling Psychophysical Data in R}. New York: Springer.

\item Murray, R.F. (2011). Classification images: A review. \textit{Journal of Vision}, 11(5), 1-25.

\item Lu, Z.L., \& Dosher, B.A. (2008). Characterizing observers using external noise and observer models: Assessing internal representations with external noise. \textit{Psychological Review}, 115(1), 44-82.

\item Prins, N. (2012). The psychometric function: The lapse rate revisited. \textit{Journal of Vision}, 12(6), 1-16.

\item Morgan, M., Dillenburger, B., Raphael, S., \& Solomon, J.A. (2012). Observers can voluntarily shift their psychometric functions without losing sensitivity. \textit{Attention, Perception, \& Psychophysics}, 74(1), 185-193.
\end{enumerate}

\end{document}
